(1) 链条不可伸长且均沿自身移动，所有链元的速度大小均为 \(v=\dot x\)，总动能是 \(K=\tfrac12\mu L\dot x^2\)。以桌面高度为重力势能零点，只有垂直悬挂的质量产生位置相关势能，其质心低于桌面 \(x/2\)，故 \(U(x)=-\tfrac12\mu gx^2\)；桌面上的链元高度不变。
(2) 拉氏量 \(\mathcal L=K-U=\tfrac12\mu L\dot x^2+\tfrac12\mu gx^2\)，Euler–Lagrange 方程是 \(\mu L\ddot x=\mu gx\)，即 \(\ddot x=(g/L)x\)。虽然驱动力仅是悬垂段重力 \(\mu gx\)，整根链条同时运动，惯性必须为 \(\mu L\)，不能把动力学错误写成 \(\mu x\ddot x=\mu gx\)。
(3) 记 \(\omega=\sqrt{g/L}\)，初值问题 \(\ddot x=\omega^2x\)、\(\dot x(0)=0\) 的解为 \(x(t)=x_0\cosh(\omega t)\)、\(\dot x(t)=x_0\omega\sinh(\omega t)\)。当最后一段将要越过桌缘时 \(x(t_*)=L\)，所以 \(t_*=\sqrt{L/g}\,\operatorname{arcosh}(L/x_0)\)；仅当 \(0<x_0<L\) 时该表达式是正的有限到达时间。
(4) 能量 \(E=\tfrac12\mu L\dot x^2-\tfrac12\mu gx^2\) 守恒，代入初态得到 \(\dot x^2=g(x^2-x_0^2)/L\)；把第三问双曲函数关系代入可逐项验证。光滑转向假设避免链段与尖锐桌缘发生冲击和局部弯曲损失；若实际桌角诱发摩擦、撞击或链条的有限弯曲刚度，需加入额外作用并重新检验能量守恒。
